% CASCADE A0 portrait poster (Beamer)
% Compile with:  latexmk -pdf cascade-poster-portrait.tex
%   (or -xelatex / -lualatex to use the Rubik heading font, if installed)
\documentclass[final]{beamer}
\usepackage[size=a0,orientation=portrait,scale=1.4]{beamerposter}
\usetheme{Cascade}
\usepackage{pgfplots}
\pgfplotsset{compat=1.17}

\title{Your poster title in up to two lines of large, readable type}
\author{Author One, Author Two, Author Three}
\institute{Department of Computer Science and Technology, University of Cambridge}
\event{Conference name \textperiodcentered{} City \textperiodcentered{} Date}
\contact{name@cl.cam.ac.uk}
\qrurl{https://www.cl.cam.ac.uk/research/cascade/}   % the QR code links here
\supporters{Supported by CASCADE's industry partners --- add supporter logos or acknowledgements here}

\begin{document}
\begin{frame}[t]
\begin{columns}[t,totalwidth=\linewidth]

% ---------------------------------------------------------------- Column 1
\begin{column}{0.31\linewidth}
  \begin{block}{Motivation}
    State the problem in one or two sentences. Why does it matter for
    computer architecture, design automation or semiconductors?
    \begin{itemize}
      \item Key challenge one
      \item Key challenge two
      \item Key challenge three
    \end{itemize}
  \end{block}

  \begin{block}{Background}
    Summarise prior work and the gap your research addresses. Keep it short:
    readers give a poster about a minute.
  \end{block}

  \begin{figure}
    \cascadeplaceholder{30cm}
    % \includegraphics[width=\linewidth]{figures/figure1}
    \caption{Describe what the reader should notice.}
  \end{figure}
  \vskip1.5em

  \begin{alertblock}{KEY TAKEAWAY}
    One sentence a passer-by should remember.
  \end{alertblock}
\end{column}

% ---------------------------------------------------------------- Column 2
\begin{column}{0.31\linewidth}
  \begin{block}{Approach}
    Explain your method, architecture or tool flow. A diagram usually works
    better than text here.
  \end{block}

  \begin{figure}
    \cascadeplaceholder{30cm}
    \caption{System overview / microarchitecture diagram.}
  \end{figure}
  \vskip2.2em

  \begin{block}{Results}
    \begin{cascadecard}
      \begin{columns}[t,totalwidth=\linewidth]
        \column{0.32\linewidth}\cascadestat{2.3$\times$}{speed-up}
        \column{0.32\linewidth}\cascadestat{$-$41\%}{energy}
        \column{0.32\linewidth}\cascadestat{$-$18\%}{silicon area}
      \end{columns}
    \end{cascadecard}
    \vskip0.6em
    Interpret the headline numbers: what was measured, against what baseline,
    on which benchmarks or silicon.
  \end{block}
\end{column}

% ---------------------------------------------------------------- Column 3
\begin{column}{0.31\linewidth}
  \begin{center}
  \begin{tikzpicture}
    \begin{axis}[
        width=\linewidth, height=26cm,
        cascade bars, bar width=1.4cm, ymax=3.6, ytick={0,1,2,3}, yticklabel style={/pgf/number format/precision=0},
        symbolic x coords={Bench A,Bench B,Bench C,Bench D}, xtick=data,
        ylabel={Speed-up}, title={\textbf{Performance vs.\ baseline}},
]
      \addplot[fill=cascadelightgrey] coordinates {(Bench A,1.0) (Bench B,1.0) (Bench C,1.0) (Bench D,1.0)};
      \addplot[fill=cascadered] coordinates {(Bench A,1.6) (Bench B,2.3) (Bench C,1.4) (Bench D,3.1)};
      \legend{Baseline, Our design}
    \end{axis}
  \end{tikzpicture}
  \end{center}
  \vskip2em

  \begin{block}{Conclusions}
    \begin{itemize}
      \item Main finding, stated plainly
      \item What it enables
      \item Limitations and open questions
    \end{itemize}
  \end{block}

  \begin{block}{Next steps}
    Where the work goes next, and how industry partners or students can get involved.
  \end{block}

  {\footnotesize\color{cascadegrey}
   \textbf{\color{cascadetext}References}\par
   [1] A.~Author et al., Paper title, \emph{Venue}, Year.\par
   [2] B.~Author et al., Paper title, \emph{Venue}, Year.\par}
\end{column}

\end{columns}
\end{frame}
\end{document}
